% GENERATED FILE. Edit canonical lessons, then run npm run generate:books.
% canonical-source-hash: fnv1a64:705c049879b555e6
% canonical-lessons: KA-C462-balkani, KA-C462-lipht, KA-C462-solleparade, KA-C462-hitar, KA-C462-kiruguttu

\chapter{In the Hotel Room}
\label{ch:ka-in-the-hotel-room}
\hlchaptermodality{462}

\begin{chapteropening}
\textbf{By the end of this chapter:} \emph{I can say a balcony, an elevator, a mosquito net, a heater and to creak.}

Complete the last lesson of chapter 462: I can say a balcony, an elevator, a mosquito net, a heater and to creak.
\end{chapteropening}

\section[\texorpdfstring{\kn{ಬಾಲ್ಕನಿ} (bālkani)}{ಬಾಲ್ಕನಿ (bālkani)}]{\kn{ಬಾಲ್ಕನಿ} (bālkani) — a balcony}

\label{lesson:KA-C462-balkani}



\begin{quote}
Before the new one: say the Kannada for a locker, then the Kannada for a refund.
\end{quote}

\subsection*{You'll want to know: \kn{ಬಾಲ್ಕನಿ}}
\textbf{\kn{ಬಾಲ್ಕನಿ}} — \emph{bālkani} — \textquotedblleft{}a balcony\textquotedblright{}.

\begin{quote}
\kn{ಕೊಠಡಿಯ} \kn{ಬಾಲ್ಕನಿಯಿಂದ} \kn{ಸಮುದ್ರ} \kn{ಕಾಣುತ್ತದೆ}.
\end{quote}

\begin{quote}
\emph{koṭhaḍiya bālkaniyinda samudra kāṇuttade.}
\end{quote}

\begin{quote}
The sea can be seen from the balcony of the room.
\end{quote}

Older houses have a \kn{ಜಗಲಿ} in front instead.

\begin{cousinweb}
\textbf{\kn{ಬಾಲ್ಕನಿ}} (\emph{bālkani}) is a loan from English \emph{balcony}.
\end{cousinweb}

\begin{grammarlens}[title={What you've built}]
One more word for getting around.
\end{grammarlens}

\subsection*{Your turn}
\begin{itemize}
\raggedright
  \item \emph{Say it:} \emph{bālkani}
  \item \emph{Say it:} \emph{bālkani}, once more
  \item \emph{Say it:} \emph{koṭhaḍiya bālkaniyinda samudra kāṇuttade.} — The sea can be seen from the balcony of the room.
  \item \emph{Recall:} say the Kannada for a locker, then the Kannada for a refund, then say \emph{bālkani} again
\end{itemize}

\begin{tcolorbox}[breakable,colback=teal!4,colframe=teal!35!black,title={Before you move on}]
What does \kn{ಬಾಲ್ಕನಿ} mean? (\textquotedblleft{}A balcony\textquotedblright{}.) Say it once more.
\end{tcolorbox}
\section[\texorpdfstring{\kn{ಲಿಫ್ಟ್} (liphṭ)}{ಲಿಫ್ಟ್ (liphṭ)}]{\kn{ಲಿಫ್ಟ್} (liphṭ) — an elevator; a lift}

\label{lesson:KA-C462-lipht}



\begin{quote}
Before the new one: say the Kannada for a refund, then the Kannada for a balcony.
\end{quote}

\subsection*{You'll want to know: \kn{ಲಿಫ್ಟ್}}
\textbf{\kn{ಲಿಫ್ಟ್}} — \emph{liphṭ} — \textquotedblleft{}an elevator; a lift\textquotedblright{}.

\begin{quote}
\kn{ಲಿಫ್ಟ್} \kn{ಕೆಲಸ} \kn{ಮಾಡದ} \kn{ಕಾರಣ} \kn{ನಾವು} \kn{ಮೆಟ್ಟಿಲು} \kn{ಹತ್ತಿದೆವು}.
\end{quote}

\begin{quote}
\emph{liphṭ kelasa māḍada kāraṇa nāvu meṭṭilu hattidevu.}
\end{quote}

\begin{quote}
Because the lift was not working we climbed the stairs.
\end{quote}

\kn{ಲಿಫ್ಟ್} \kn{ಕೊಡು} is also to give someone a ride.

\begin{cousinweb}
\textbf{\kn{ಲಿಫ್ಟ್}} (\emph{lipṭ}) is a loan from English \emph{lift}.
\end{cousinweb}

\begin{grammarlens}[title={What you've built}]
One more word for getting around.
\end{grammarlens}

\subsection*{Your turn}
\begin{itemize}
\raggedright
  \item \emph{Say it:} \emph{liphṭ}
  \item \emph{Say it:} \emph{liphṭ}, once more
  \item \emph{Say it:} \emph{liphṭ kelasa māḍada kāraṇa nāvu meṭṭilu hattidevu.} — Because the lift was not working we climbed the stairs.
  \item \emph{Recall:} say the Kannada for a refund, then the Kannada for a balcony, then say \emph{liphṭ} again
\end{itemize}

\begin{tcolorbox}[breakable,colback=teal!4,colframe=teal!35!black,title={Before you move on}]
What does \kn{ಲಿಫ್ಟ್} mean? (\textquotedblleft{}An elevator; a lift\textquotedblright{}.) Say it once more.
\end{tcolorbox}
\section[\texorpdfstring{\kn{ಸೊಳ್ಳೆಪರದೆ} (soḷḷeparade)}{ಸೊಳ್ಳೆಪರದೆ (soḷḷeparade)}]{\kn{ಸೊಳ್ಳೆಪರದೆ} (soḷḷeparade) — a mosquito net}

\label{lesson:KA-C462-solleparade}



\begin{quote}
Before the new one: say the Kannada for a balcony, then the Kannada for an elevator.
\end{quote}

\subsection*{You'll want to know: \kn{ಸೊಳ್ಳೆಪರದೆ}}
\textbf{\kn{ಸೊಳ್ಳೆಪರದೆ}} — \emph{soḷḷeparade} — \textquotedblleft{}a mosquito net\textquotedblright{}.

\begin{quote}
\kn{ಮಂಚದ} \kn{ಮೇಲೆ} \kn{ಒಂದು} \kn{ಸೊಳ್ಳೆಪರದೆ} \kn{ಇತ್ತು}.
\end{quote}

\begin{quote}
\emph{mañcada mēle ondu soḷḷeparade ittu.}
\end{quote}

\begin{quote}
There was a mosquito net over the bed.
\end{quote}

\kn{ಸೊಳ್ಳೆ}, a mosquito, plus \kn{ಪರದೆ}, a curtain.

\begin{cousinweb}
\textbf{\kn{ಸೊಳ್ಳೆಪರದೆ}} (\emph{soḷḷeparade}) is a Kannada compound of \kn{ಸೊಳ್ಳೆ}, \textquotedblleft{}a mosquito\textquotedblright{}, and \kn{ಪರದೆ}, \textquotedblleft{}a curtain\textquotedblright{}, which is usually traced to Persian \emph{parda}.
\end{cousinweb}

\begin{grammarlens}[title={What you've built}]
One more word for getting around.
\end{grammarlens}

\subsection*{Your turn}
\begin{itemize}
\raggedright
  \item \emph{Say it:} \emph{soḷḷeparade}
  \item \emph{Say it:} \emph{soḷḷeparade}, once more
  \item \emph{Say it:} \emph{mañcada mēle ondu soḷḷeparade ittu.} — There was a mosquito net over the bed.
  \item \emph{Recall:} say the Kannada for a balcony, then the Kannada for an elevator, then say \emph{soḷḷeparade} again
\end{itemize}

\begin{tcolorbox}[breakable,colback=teal!4,colframe=teal!35!black,title={Before you move on}]
What does \kn{ಸೊಳ್ಳೆಪರದೆ} mean? (\textquotedblleft{}A mosquito net\textquotedblright{}.) Say it once more.
\end{tcolorbox}
\section[\texorpdfstring{\kn{ಹೀಟರ್} (hīṭar)}{ಹೀಟರ್ (hīṭar)}]{\kn{ಹೀಟರ್} (hīṭar) — a heater}

\label{lesson:KA-C462-hitar}



\begin{quote}
Before the new one: say the Kannada for an elevator, then the Kannada for a mosquito net.
\end{quote}

\subsection*{You'll want to know: \kn{ಹೀಟರ್}}
\textbf{\kn{ಹೀಟರ್}} — \emph{hīṭar} — \textquotedblleft{}a heater\textquotedblright{}.

\begin{quote}
\kn{ಕೊಠಡಿಯಲ್ಲಿ} \kn{ಹೀಟರ್} \kn{ಇದ್ದರೂ} \kn{ಚಳಿ} \kn{ಆಗುತ್ತಿತ್ತು}.
\end{quote}

\begin{quote}
\emph{koṭhaḍiyalli hīṭar iddarū caḷi āguttittu.}
\end{quote}

\begin{quote}
Even with a heater in the room it felt cold.
\end{quote}

\kn{ಹೀಟರ್} \kn{ಹಾಕು}, to put the heater on.

\begin{cousinweb}
\textbf{\kn{ಹೀಟರ್}} (\emph{hīṭar}) is a loan from English \emph{heater}.
\end{cousinweb}

\begin{grammarlens}[title={What you've built}]
One more word for getting around.
\end{grammarlens}

\subsection*{Your turn}
\begin{itemize}
\raggedright
  \item \emph{Say it:} \emph{hīṭar}
  \item \emph{Say it:} \emph{hīṭar}, once more
  \item \emph{Say it:} \emph{koṭhaḍiyalli hīṭar iddarū caḷi āguttittu.} — Even with a heater in the room it felt cold.
  \item \emph{Recall:} say the Kannada for an elevator, then the Kannada for a mosquito net, then say \emph{hīṭar} again
\end{itemize}

\begin{tcolorbox}[breakable,colback=teal!4,colframe=teal!35!black,title={Before you move on}]
What does \kn{ಹೀಟರ್} mean? (\textquotedblleft{}A heater\textquotedblright{}.) Say it once more.
\end{tcolorbox}
\section[\texorpdfstring{\kn{ಕಿರುಗುಟ್ಟು} (kiruguṭṭu)}{ಕಿರುಗುಟ್ಟು (kiruguṭṭu)}]{\kn{ಕಿರುಗುಟ್ಟು} (kiruguṭṭu) — to creak; to squeak}

\label{lesson:KA-C462-kiruguttu}



\begin{quote}
Before the new one: say the Kannada for a mosquito net, then the Kannada for a heater.
\end{quote}

\subsection*{You'll want to know: \kn{ಕಿರುಗುಟ್ಟು}}
\textbf{\kn{ಕಿರುಗುಟ್ಟು}} — \emph{kiruguṭṭu} — \textquotedblleft{}to creak; to squeak\textquotedblright{}.

\begin{quote}
\kn{ರಾತ್ರಿ} \kn{ಕೊಠಡಿಯ} \kn{ಹಳೆಯ} \kn{ಬಾಗಿಲು} \kn{ಕಿರುಗುಟ್ಟಿತು}.
\end{quote}

\begin{quote}
\emph{rātri koṭhaḍiya haḷeya bāgilu kiruguṭṭitu.}
\end{quote}

\begin{quote}
At night the old door of the room creaked.
\end{quote}

Past (it …): \kn{ಕಿರುಗುಟ್ಟಿತು}, \emph{kiruguṭṭitu}.

Of doors, brakes, sandals and birds: a thin, sharp sound.

\begin{cousinweb}
\textbf{\kn{ಕಿರುಗುಟ್ಟು}} (\emph{kiruguṭṭu}) is native Dravidian: DEDR 1590 lists Kannada \emph{kiṟuguṭṭu}, \emph{kirguṭṭu} among words for crying out and shrill sounds.
\end{cousinweb}

\begin{grammarlens}[title={What you've built}]
One more word for getting around.
\end{grammarlens}

\subsection*{Your turn}
\begin{itemize}
\raggedright
  \item \emph{Say it:} \emph{kiruguṭṭu}
  \item \emph{Say it:} \emph{kiruguṭṭu}, once more
  \item \emph{Say it:} \emph{rātri koṭhaḍiya haḷeya bāgilu kiruguṭṭitu.} — At night the old door of the room creaked.
  \item \emph{Recall:} say the Kannada for a mosquito net, then the Kannada for a heater, then say \emph{kiruguṭṭu} again
\end{itemize}

\begin{tcolorbox}[breakable,colback=teal!4,colframe=teal!35!black,title={Before you move on}]
What does \kn{ಕಿರುಗುಟ್ಟು} mean? (\textquotedblleft{}To creak; to squeak\textquotedblright{}.) Say it once more.
\end{tcolorbox}
